% !TEX root = chapter-standalone.tex

\section{Comonads in bicategories}
\label{sec:bicats-comonads}

In \cref{sec:bicats-monads} we defined monads in bicategories and saw that they unify monads on categories, closure operators, monoid objects, and other structures.
We now consider the dual notion: \emph{comonads} in bicategories.
Just as a comonad on a category (\cref{def:comonad}) is obtained from the definition of a monad by reversing all the 2-cells, a comonad in a bicategory is obtained from \cref{def:bicat-monad} by the same reversal.

\subsection{Motivation}

Recall that a \SY{comonad} on a category~$\CatC$ (\cref{def:comonad}) consists of an endofunctor~$\comonW \colon \CatC \fto \CatC$ together with natural transformations for \emph{comultiplication} (duplication)~$\comondu \colon \comonW \nto \comonWW$ and \emph{counit} (extraction)~$\comonex \colon \comonW \nto \funidC$, satisfying coassociativity and counitality.
This is the ``arrow-reversed'' version of a monad: the multiplication~$\moncomp \colon \monA \mthen \monA \nto \monA$ becomes a comultiplication~$\comondu \colon \comonW \nto \comonW \mthen \comonW$, and the unit~$\monunit \colon \catidof{\Obja} \nto \monA$ becomes a counit~$\comonex \colon \comonW \nto \catidof{\Obja}$.

While monads capture the idea of ``wrapping'' values with additional structure, comonads capture the dual idea of ``extracting'' a value from a structured context and ``duplicating'' that context.
This duality generalizes to any bicategory: a comonad in a bicategory is an endo-1-cell equipped with a comultiplication and counit satisfying the reversed axioms.

\subsection{The definition}

\begin{ctdefinition}[Comonad in a bicategory]
    \SYNDEF{comonad in a bicategory}
    \label{def:bicat-comonad}
    Let~$\CatB$ be a bicategory and let~$\Obja \setin \Ob(\CatB)$.
    A \maindef{comonad} on~$\Obja$ in~$\CatB$ is:

    \constit

    \begin{enumerate}
        \item A 1-cell~$\comonW \colon \Obja \mto \Obja$ (an endo-1-cell);
        \item A 2-cell~$\comondu \colon \comonW \nto \comonW \mthen \comonW$, called the \emph{comultiplication} (or \emph{duplication});
        \item A 2-cell~$\comonex \colon \comonW \nto \catidof{\Obja}$, called the \emph{counit} (or \emph{extraction}).
    \end{enumerate}

    \condit

    \begin{enumerate}
        \item \emph{Coassociativity:}
              The following diagram of 2-cells commutes:
              \begin{equation}\label{eq:bicat-comonad-coassoc}
                  \begin{tikzcd}[column sep=large, row sep=large]
                      {\comonW} \ar[rr, "{\comondu}"] \ar[d, "{\comondu}"']   & \quad &   {\comonW \mthen \comonW} \ar[d, "{\catidof{\comonW} \nthenhor \comondu}"] \\
                      {\comonW \mthen \comonW} \ar[r, "{\comondu \nthenhor \catidof{\comonW}}"']  &    {(\comonW \mthen \comonW) \mthen \comonW} \ar[r, "{\associator}"']   & {\comonW \mthen (\comonW \mthen \comonW)} 
                  \end{tikzcd}
              \end{equation}

        \item \emph{Left counitality:}
              The following diagram of 2-cells commutes:
              \begin{equation}\label{eq:bicat-comonad-left-counit}
                  \begin{tikzcd}[column sep=large]
                      {\comonW}
                      \ar[r, "{\comondu}"]
                      \ar[dr, "{\leftunitor_\comonW^{-1}}"']
                      &
                      {\comonW \mthen \comonW}
                      \ar[d, "{\comonex \nthenhor \catidof{\comonW}}"] \\
                      &
                      {\catidof{\Obja} \mthen \comonW}
                  \end{tikzcd}
              \end{equation}

        \item \emph{Right counitality:}
              The following diagram of 2-cells commutes:
              \begin{equation}\label{eq:bicat-comonad-right-counit}
                  \begin{tikzcd}[column sep=large]
                      {\comonW}
                      \ar[r, "{\comondu}"]
                      \ar[dr, "{\rightunitor_\comonW^{-1}}"']
                      &
                      {\comonW \mthen \comonW}
                      \ar[d, "{\catidof{\comonW} \nthenhor \comonex}"] \\
                      &
                      {\comonW \mthen \catidof{\Obja}}
                  \end{tikzcd}
              \end{equation}
    \end{enumerate}

\end{ctdefinition}

\begin{remark}
    \label{rem:bicat-comonad-duality}
    Comparing \cref{def:bicat-comonad} with \cref{def:bicat-monad}, one sees that the definition of a comonad is obtained by reversing the direction of the 2-cells:
    \begin{itemize}
        \item the multiplication~$\moncomp \colon \monA \mthen \monA \nto \monA$ of a monad becomes the comultiplication~$\comondu \colon \comonW \nto \comonW \mthen \comonW$ of a comonad;
        \item the unit~$\monunit \colon \catidof{\Obja} \nto \monA$ becomes the counit~$\comonex \colon \comonW \nto \catidof{\Obja}$.
    \end{itemize}
    More precisely, a comonad on~$\Obja$ in~$\CatB$ is the same thing as a monad on~$\Obja$ in the bicategory~$\CatB^{\text{co}}$, which is obtained from~$\CatB$ by reversing all 2-cells (i.e., replacing each hom-category~$\CatB(\Obja, \Objb)$ by its opposite~$\CatB(\Obja, \Objb)\catop$) while keeping the 1-cells as they are.
\end{remark}

\begin{remark}
    \label{rem:bicat-comonad-comonoid}
    By analogy with \cref{rem:bicat-monad-monoid-analogy}, a comonad on~$\Obja$ in~$\CatB$ is a \emph{comonoid} in the monoidal category~$\CatB(\Obja, \Obja)$, where the monoidal product is horizontal composition and the monoidal unit is~$\catidof{\Obja}$.
\end{remark}

\begin{remark}
    \label{rem:bicat-comonad-strict}
    In a strict 2-category, the associators and unitors are identities, so the diagrams simplify.
    In the strict 2-category~$\Category$, a comonad on an object~$\CatC$ in the sense of \cref{def:bicat-comonad} consists of an endofunctor~$\comonW \colon \CatC \fto \CatC$, natural transformations~$\comondu \colon \comonW \nto \comonWW$ and~$\comonex \colon \comonW \nto \funidC$, satisfying coassociativity and counitality --- this is exactly \cref{def:comonad}.
\end{remark}

\subsection{Comonads from adjunctions}

Just as every adjunction gives rise to a monad (\cref{prop:bicat-adj-gives-monad}), every adjunction also gives rise to a comonad --- but on the \emph{other} object.

\begin{proposition}
    \label{prop:bicat-adj-gives-comonad}
    Let~$\CatB$ be a bicategory and let~$\funl \adjunction \funr$ be an adjunction in~$\CatB$ as in \cref{def:bicat-adjunction}, with~$\funl \colon \Obja \mto \Objb$ and~$\funr \colon \Objb \mto \Obja$.
    Then the endo-1-cell~$\comonW = \funr \mthen \funl \colon \Objb \mto \Objb$ carries the structure of a comonad on~$\Objb$, with:
    \begin{itemize}
        \item \emph{Counit:} the counit~$\equivcounit \colon \funr \mthen \funl \nto \catidof{\Objb}$ of the adjunction;
        \item \emph{Comultiplication:} the composite 2-cell
              \begin{equation}\label{eq:bicat-comonad-from-adj-comult}
                  \comonW = \funr \mthen \funl \overset{\catidof{\funr} \nthenhor \leftunitor_\funl^{-1}}{\ntolong} \funr \mthen (\catidof{\Obja} \mthen \funl) \overset{\catidof{\funr} \nthenhor (\equivunit \nthenhor \catidof{\funl})}{\ntolong} \funr \mthen ((\funl \mthen \funr) \mthen \funl) \overset{\associator^{-1}}{\ntolong} (\funr \mthen \funl) \mthen (\funr \mthen \funl) = \comonW \mthen \comonW,
              \end{equation}
              where we have suppressed some associators for readability.
    \end{itemize}
    The comonad axioms follow from the triangle identities of the adjunction.
\end{proposition}

\begin{remark}
    Note the asymmetry: from an adjunction~$\funl \adjunction \funr$ between~$\Obja$ and~$\Objb$, we obtain a \emph{monad}~$\funl \mthen \funr$ on~$\Obja$ and a \emph{comonad}~$\funr \mthen \funl$ on~$\Objb$.
    The monad uses the counit to define its multiplication, while the comonad uses the unit to define its comultiplication.
    In the strict 2-category~$\Category$, this recovers the classical fact that any adjunction~$\funl \adjunction \funr$ yields both a monad on the domain and a comonad on the codomain.
\end{remark}

\subsection{Examples}

\begin{example}[Comonads in \Category]
    \label{exa:bicat-comonad-cat}
    A comonad in the strict 2-category~$\Category$ is precisely a comonad on a category in the classical sense of \cref{def:comonad}: an endofunctor equipped with a comultiplication and counit satisfying coassociativity and counitality.
    Examples include the product (environment) comonad, the stream comonad, and the store comonad, all of which are discussed in detail in the section on comonads for ordinary categories.
\end{example}

\begin{example}[Comonads in a locally posetal bicategory: interior operators]
    \label{exa:bicat-comonad-interior}
    Consider the locally posetal bicategory where the objects are posets, the 1-cells are monotone maps, and the 2-cells are given by the pointwise ordering.
    A comonad on a poset~$\posgenA$ in this bicategory is a monotone map~$\comonW \colon \posgenA \mto \posgenA$ together with:
    \begin{itemize}
        \item a counit:~$\comonW \leq \catidof{\posgenA}$, meaning~$\comonW(x) \leq x$ for all~$x$ (\emph{deflation});
        \item a comultiplication:~$\comonW \leq \comonW \mthen \comonW$, meaning~$\comonW(x) \leq \comonW(\comonW(x))$ for all~$x$ (\emph{idempotency}, when combined with deflation).
    \end{itemize}
    Since the hom-categories are posets, the coassociativity and counitality conditions are automatic.
    This is precisely the definition of an \emph{interior operator} (also called a \emph{kernel operator}) on~$\posgenA$.
    Thus: \emph{interior operators are comonads in a locally posetal bicategory}.

    This is dual to \cref{exa:bicat-monad-closure}: closure operators are monads, and interior operators are comonads.
    Given a Galois connection~$\funl \adjunction \funr$ between posets~$\posgenA$ and~$\posgenB$ (see \cref{exa:bicat-adj-galois}), the composite~$\funl \mthen \funr$ is a closure operator on~$\posgenA$ (a monad), while the composite~$\funr \mthen \funl$ is an interior operator on~$\posgenB$ (a comonad), in accordance with \cref{prop:bicat-adj-gives-comonad}.
\end{example}

\begin{example}[Comonads in a one-object bicategory: comonoids in monoidal categories]
    \label{exa:bicat-comonad-monoidal}
    Recall from \cref{exa:bicat-monoidal} that a monoidal category~$\tup{\CatC, \mtimescat, \idmoncat}$ corresponds to a one-object bicategory~$\mathbf{\Sigma}\CatC$.
    A comonad in~$\mathbf{\Sigma}\CatC$ is an object~$W$ of~$\CatC$ equipped with morphisms
    \begin{equation}
        \comondu \colon W \mto W \mtimescatob W \qqand \comonex \colon W \mto \idmoncat
    \end{equation}
    satisfying coassociativity and counitality.
    This is precisely a \emph{comonoid object} in the monoidal category~$\CatC$.

    For instance, in the monoidal category~$(\Set, \cartprod, \singleton)$ of sets with cartesian product, \emph{every} set carries a unique comonoid structure: the comultiplication is the diagonal~$\ela \mapsto (\ela, \ela)$ and the counit is the unique map to the singleton.
    This reflects the fact that in~$\Set$ one can always ``copy'' and ``discard'' data --- a property that fails in other monoidal categories, such as that of vector spaces with the tensor product.
\end{example}

\begin{example}[Comonads in the bicategory of Mealy machines]
    \label{exa:bicat-comonad-mealy}
    A comonad on a set~$\setA$ in the bicategory~$\operatorname{\textbf{Mealy}}$ (\cref{exa:bicat-mealy}) is a Mealy machine~$\comonW \colon \setA \mto \setA$ together with simulations
    \begin{equation}
        \comondu \colon \comonW \nto \comonW \mthen \comonW \qqand \comonex \colon \comonW \nto \catidof{\setA},
    \end{equation}
    satisfying coassociativity and counitality.
    The extraction simulation~$\comonex$ maps the state space of~$\comonW$ into the trivial state space of the identity machine, expressing that the machine's behavior can be ``observed'' as the identity.
    The duplication simulation~$\comondu$ maps the state space of~$\comonW$ into the product state space of two copies composed in series, expressing that the machine's state can be coherently ``split'' between two copies of itself.

    Intuitively, while a monad structure on a Mealy machine says that iterating the machine can be collapsed into a single run (\cref{exa:bicat-monad-mealy}), a comonad structure says that a single run can be coherently split into two successive runs.
\end{example}

\subsection{Summary}

The following table extends the summary from \cref{sec:bicats-monads} by adding the comonad column.

\begin{center}
    \begin{tabular}{llll}
        \textbf{Bicategory} & \textbf{Adjunction} & \textbf{Monad} & \textbf{Comonad} \\
        \hline
        $\Category$ & adjunction & monad & comonad \\
        locally posetal & Galois connection & closure operator & interior operator \\
        one-object ($\mathbf{\Sigma}\CatC$) & dual pair & monoid in~$\CatC$ & comonoid in~$\CatC$ \\
        $\operatorname{\textbf{Mealy}}$ & machines inverse up to sim. & collapsible iteration & splittable run \\
    \end{tabular}
\end{center}

The general pattern is: whenever an adjunction~$\funl \adjunction \funr$ exists in a bicategory, the composite~$\funl \mthen \funr$ carries a monad structure and the composite~$\funr \mthen \funl$ carries a comonad structure.
Monads and comonads are thus two sides of the same coin, each arising from one of the two composites of an adjunction.
